\chapter{The Closing Statement}
\label{ch:synthesis}

\begin{quotation}
\noindent\textit{Synopsis.} This closing chapter integrates the
monograph's concepts into a single narrative: distinguishability explains
what survives; admissibility explains what remains possible; repair
economics explains why a distinction was or was not maintained; budget
relays explain how losses accumulate across time and institutions; the
fidelity continuum unifies preservation and reconstitution; and Composition
Failure explains why locally rational decisions can still produce globally
insufficient outcomes. The chapter closes by naming what remains
unresolved, above all the MEM|8 dependency, confronting directly the
strongest available objection to the whole enterprise -- that its central
theorem may be little more than task-unavailability restated in heavier
notation -- and restates the central claim
that continuation is constrained not by malice or error but by the
unavoidable limits of finite repair budgets operating in an open future.
\end{quotation}

\section{Restating the Arc}

Distinguishability (\cref{ch:distinctions}) established what survives a
projection. Admissibility (\cref{ch:admissibility}) established what
continuations remain reachable given what survives. Repair economics
(\cref{ch:repair-economics}) established why a given distinction was or
was not maintained, and, via \cref{cor:ambiguity-limiting-case},
established that this monograph could not simply inherit its central
claim from \emph{Conservation of Ambiguity}. The budget relay
(\cref{ch:budget-relay}) formalized how discard compounds across
uncoordinated stages, and the Composition Failure Theorem
(\cref{thm:composition-failure}) proved the diachronic analogue the
Redistribution Theorem does not cover. The fidelity continuum
(\cref{ch:fidelity-continuum}) collapsed preservation and
reconstitution into a single variable, conditional on the unresolved
question raised in Chapter~\ref{ch:mem8}.

Part~\ref{part:worked-domains} then did more than illustrate this arc.
Chess and maps showed the arc's positive case: sufficiency genuinely
achieved, by augmentation or by declared discard, when hidden state is
closed. Painting showed its negative case, and showed that the negative
case is not a defect: ambiguity a task does not need to resolve is not a
gap to be regretted. Proofs supplied the arc's sharpest internal check,
\cref{prop:homogeneous-chain}, which deserves to sit alongside
\cref{thm:composition-failure} in this summary rather than beneath it:
where the Composition Failure Theorem shows how differing local tasks
compose into a gap no stage intended, \cref{prop:homogeneous-chain}
shows the exact condition under which they do not -- a single task,
adequately funded, held constant across every stage. The two results are
converses of each other, and together they say more than either does
alone: composition failure is not a property of chains in general, it is
a property of chains whose stages do not agree, deliberately or by
accident, on what matters. Contracts, science, memory, and pipelines
each instantiated this same pair of results in a different institutional
setting, and the pipeline chapter's numerically explicit witness
(\cref{prop:pipeline-witness}) closed the loop back to
Appendix~\ref{app:proof}'s abstract construction, confirming the two are
the same argument rather than merely similar ones.

\section{The Closing Statement}

\begin{conjecture}[Closing Statement]
\label{conj:closing}
Continuation depends on the substrate available to a consumer, whether
preserved with high fidelity or regenerated from traces at low fidelity;
in either case the substrate is bounded by the distinctions its
production budget could afford to maintain. Because relays compose
without guaranteed coordination (\cref{def:budget-relay}) and because
future task spaces remain open (\cref{prop:task-unavailability}), no
finite chain of substrates can be guaranteed sufficient for arbitrary
future continuation -- not because information is destroyed, but
because no stage can afford, or is positioned, to preserve everything a
task it cannot yet foresee will require.
\end{conjecture}

This statement is labeled a conjecture rather than a theorem
deliberately: it synthesizes \cref{thm:composition-failure} together
with the fidelity continuum of Chapter~\ref{ch:fidelity-continuum}, but
the synthesis itself, as a single unified statement, has not been given
an independent proof beyond what its components already establish.

\section{What Remains Open}

Appendix~\ref{app:open-questions} consolidates every conjecture and open
question raised in the body in one register, organized by which part of
the framework each bears on rather than by chapter, and this section
does not attempt to duplicate it. Three items are worth naming here
because of their scope rather than their difficulty. The MEM|8 fork
(\cref{oq:mem8-fork}) is the largest: it bears specifically on whether a
consumer's perceptual access to a given rendering can itself reach
fidelity $1$, not on whether an already-open task's ambiguity can be
resolved by perceptual directness alone -- \cref{rem:fork-scope} states
this boundary precisely, and it is worth a reader's attention, since
conflating the two was an error this monograph made and corrected during
its own drafting. The tightness of \cref{cor:no-free-lunch}'s growth
bound is conjectural: this monograph shows cumulative discard does not
shrink, not that discard \emph{risk} grows at any particular rate.
Whether a general coordination protocol across budget relays could
bound, though not eliminate, \cref{thm:composition-failure}'s failure
mode is left as a direction for future work; \cref{prop:pooling-restores}
shows one specific, more radical alternative to coordination succeeding
in one instance, which is evidence that the broader question is not
hopeless, not an answer to it.

\section{What This Monograph Does Not Claim}

This monograph does not claim that all cognition is reconstructive; it
restricts that claim, via \cref{ch:mem8}'s narrow reading, to cases
independently supported by existing literature, and Chapter~\ref{ch:memory}
found, on inspection, that none of its own conclusions actually required
the fork to be resolved either way. It does not claim that
\emph{Conservation of Ambiguity}'s Redistribution Theorem is superseded;
\cref{thm:redistribution} is imported and used, not revised, and
Appendix~\ref{app:comparison} shows concretely, with a specific
allocation of resources that succeeds under pooling where the separated
relay fails, that the two theorems answer different questions rather
than one subsuming the other. It does not claim to have resolved the
preservation/reconstitution question definitively, since that resolution
remains contingent on \cref{oq:mem8-fork}. And it does not claim that
the closing statement above is the only way to unify its results; it
claims only that this particular unification is not, on inspection,
derivable from any single component alone, which is the most a
conjecture honestly labeled as such should claim.

\section{The Strongest Objection}

The most serious challenge to this monograph is not a technical one, and
it deserves a direct answer rather than a deflection into further
formalism. The objection: \cref{thm:composition-failure}, once stated
precisely, is a fairly direct consequence of
\cref{prop:task-unavailability} -- no substrate can anticipate an
unforeseeable future task -- dressed in the additional vocabulary of
budgets, stages, and discard sets. If that is right, has this monograph
built an elaborate apparatus to reprove something close to common sense:
that you cannot plan for what you do not know?

The honest answer has three parts, and none of them is a denial that
the objection identifies something real. First, the \emph{existence}
half of \cref{thm:composition-failure} is indeed a close relative of
\cref{prop:task-unavailability}, and this monograph has not claimed
otherwise; Chapter~\ref{ch:composition-failure}'s own proof sketch is
built directly on top of it. Second, and this is where the objection
proves too much if pressed further, the theorem's genuine content is not
the existence claim but two things the existence claim alone does not
supply: the exact mechanism by which failure occurs (budget exhaustion
after funding a stage's own required distinctions, not some
unspecified misfortune), robust to the choice of local discard procedure
(\cref{lem:exhaustion}), and the precise converse identifying when
failure provably does not occur (\cref{prop:homogeneous-chain}), which
is not a consequence of \cref{prop:task-unavailability} at all and could
not have been stated without the budget-relay apparatus this monograph
built to state it. A theorem that only says ``the unforeseeable cannot
be foreseen'' does not tell you that a shared, adequately funded task
throughout a chain is exactly what defends against it; this monograph's
apparatus does. Third, and this is the part worth being fully candid
about: whether that additional content justifies the apparatus's
formal weight, as opposed to a shorter and less machinery-heavy
statement of the same two facts, is a matter of taste this monograph
cannot settle for a skeptical reader by further argument. The worked
domains of Part~\ref{part:worked-domains} are this monograph's actual
case for the apparatus's usefulness -- chess, painting, proofs, and
pipelines each get a sharper account of their own structure from having
the vocabulary available -- and a reader unconvinced by those chapters
will not be talked into the apparatus's worth by this section restating
that it exists.

\section{Future Research Directions}

Appendix~\ref{app:open-questions} is the complete register; this section
does not duplicate it, only names the three items there most likely to
reward attention first. The tightness of \cref{cor:no-free-lunch}'s
growth bound is the most tractable: it asks a concrete quantitative
question about a construction already fully specified in
Appendix~\ref{app:proof}. The measure of a persistence set is the most
promising: \cref{rem:witnesses-persistence-measure}'s reconstruction
redundancy $\rho_W(r)$, though unverified, is concrete enough to check
directly against this monograph's root-recovery apparatus
(\cref{prop:root-recovery}) without requiring new machinery. The MEM|8
fork is the most consequential and the least tractable by this
monograph's methods: resolving it requires empirical and philosophical
work well outside what a framework built from worked examples and
formal proof can settle on its own, and this monograph's contribution is
limited to having shown precisely which of its own results would move,
and which would not, once someone else resolves it.
